\documentclass[11pt,a4paper]{article}

% ---------------------------------------------------------------------------
% NeuroShield research paper. Self-contained: uses only widely available
% packages, so it compiles on Overleaf or any standard TeX Live install with
% no .bib file (references are a manual thebibliography).
% ---------------------------------------------------------------------------

\usepackage[utf8]{inputenc}
\usepackage[T1]{fontenc}
\usepackage[a4paper,margin=1in]{geometry}
\usepackage{amsmath,amssymb}
\usepackage{booktabs}
\usepackage{array}
\usepackage[table]{xcolor}
\usepackage{graphicx}
\usepackage{caption}
\usepackage{enumitem}
\usepackage{titlesec}
\usepackage[hidelinks]{hyperref}

\definecolor{barfill}{HTML}{3B5B92}
\definecolor{rowalt}{HTML}{F2F4F7}

\titleformat{\section}{\normalfont\large\bfseries}{\thesection}{0.6em}{}
\titleformat{\subsection}{\normalfont\normalsize\bfseries}{\thesubsection}{0.6em}{}
\setlength{\parskip}{0.35em}
\setlength{\parindent}{0pt}

\newcommand{\kpi}[1]{\textbf{#1}}

\title{\vspace{-1.2cm}\textbf{NeuroShield: A Personalized, Multi-Signal System for
Honest Real-Time Stress Detection}}
\author{Manika Jhaveri}
\date{\today}

\begin{document}
\maketitle

% ===========================================================================
\begin{abstract}
\noindent
I present NeuroShield, a system that estimates a person's physiological stress in real time from
low-cost wearable-grade sensors, and does so \emph{relative to that individual's own resting
baseline} rather than a population average. The motivating observation is that absolute
physiological values are ambiguous across people: two individuals can both be fully at rest at
$58$ and $82$ beats per minute, so a model fed raw values expends capacity learning \emph{who} a
person is instead of \emph{how aroused} they are. NeuroShield calibrates to each user during a short
quiet period and expresses every subsequent reading as a deviation from that baseline. Every
$60$~seconds it derives $36$ features from four physiological channels---photoplethysmography (PPG),
electrodermal activity (EDA), skin temperature, and accelerometry---and feeds them to a two-headed
gradient-boosted model: a calibrated head that emits a $0$--$100$ graded stress index, and a
four-class head that separates stress arousal from positive arousal such as amusement. Under
grouped leave-one-subject-out (LOSO) evaluation---training on some subjects and testing only on
subjects never seen during training---the graded-stress head attains a balanced accuracy of
\kpi{0.919}, exceeding the strongest comparable wrist-only result reported in the literature
($0.874$). Beyond raw accuracy, the system is designed for \emph{epistemic honesty}: it abstains
from prediction when motion or poor signal quality render a reading untrustworthy, it surfaces its
own uncertainty (a $\pm 0.036$ standard error and a per-subject accuracy range of $0.50$--$1.00$),
and it is structurally prevented from issuing any clinical or diagnostic claim. I describe the
research, problem, objectives, plan, methodology, datasets, and results behind the system, and I
report a negative result establishing that $0.919$ is the practical ceiling for the available data.
\end{abstract}

\vspace{0.3em}
\noindent\textbf{Keywords:} affective computing, stress detection, wearable sensors, heart-rate
variability, electrodermal activity, personalization, gradient boosting, model calibration,
leave-one-subject-out validation.

% ===========================================================================
\section{Introduction}

Chronic stress is among the most prevalent and costly health burdens, yet the instruments used to
track it are weak. The dominant tool---self-report---is unreliable: people are poor at noticing
their own stress, and by the time they consciously register it they have often been under load for
hours. The autonomic nervous system, however, responds first, and it responds \emph{measurably}.
Psychological stress activates the sympathetic branch of the autonomic nervous system, producing a
consistent peripheral signature: elevated heart rate, suppressed heart-rate variability, increased
sweat-gland activity, and cooling of the skin as blood is redirected toward the core. Each of these
is observable with inexpensive sensors.

The difficulty is that turning those signals into a trustworthy estimate of stress is far harder
than it first appears, for three reasons. First, physiological signals vary enormously between
individuals, so a fixed rule (``heart rate above $90 \Rightarrow$ stress'') is meaningless across a
population. Second, wearable sensors on a moving body produce unreliable data much of the time, and
a model that outputs a confident number during that period is not merely useless but actively
misleading. Third, the commercial temptation is to claim detection of ``anxiety'' or ``burnout,''
which the underlying signals physically cannot support.

NeuroShield addresses all three. My contributions are:

\begin{enumerate}[leftmargin=1.4em,itemsep=0.15em]
  \item A \textbf{personalization method} that expresses every feature relative to each user's own
        calibrated baseline, which I show is the single largest driver of cross-subject accuracy.
  \item A \textbf{two-headed model} that produces a calibrated, graded ($0$--$100$) stress index and
        additionally distinguishes stress arousal from positive arousal.
  \item An \textbf{abstention mechanism} that makes refusal to predict a first-class output.
  \item An \textbf{honest evaluation}: LOSO throughout, a reported error bar, held-out datasets, and
        a documented negative result showing that further tuning does not improve on $0.919$.
  \item An \textbf{end-to-end streaming system} in which hardware acts only as a data source, so the
        entire pipeline is validated before any device is attached.
\end{enumerate}

% ===========================================================================
\section{Background and Related Work}

I grounded the design in a review of physiological stress sensing, from which three findings shaped
every subsequent decision.

\textbf{Stress has a measurable, multi-channel signature.} The sympathetic response elevates heart
rate and reduces beat-to-beat variability, because the parasympathetic (vagal) ``brake'' on the
heart is released. It increases eccrine sweat-gland activity, raising skin conductance. And it
drives peripheral vasoconstriction, cooling the skin at the extremities. A single channel is
ambiguous; the combination is informative. Electrodermal activity is of particular value because
eccrine sweat glands are innervated \emph{solely} by the sympathetic branch, with no antagonistic
parasympathetic input~\cite{greco2016cvxeda}, making EDA the cleanest available proxy for arousal.

\textbf{Reported accuracies are frequently inflated by weak validation.} Many studies report
$93$--$99\%$ accuracy, but these commonly rely on $k$-fold or intra-subject cross-validation, which
allows the same individual to appear in both training and test partitions. A model can then memorize
person-specific idiosyncrasies rather than learn a generalizable mapping. The honest protocol is
\emph{leave-one-subject-out} (LOSO), in which the test subject is entirely unseen during training.
Under LOSO, the strongest comparable wrist-only result I identified was $0.874$~\cite{siirtola2019}.
I adopted LOSO as the sole evaluation metric.

\textbf{Personalization is the dominant lever, and naive data pooling backfires.} The wearable
affect literature, beginning with the WESAD benchmark~\cite{schmidt2018wesad}, repeatedly identifies
subject-relative normalization as the largest source of cross-subject gain. I further found---and
later empirically confirmed---that merely concatenating datasets to obtain ``more data'' can
\emph{reduce} accuracy, because it folds cross-protocol distribution shift into the model without
domain adaptation.

% ===========================================================================
\section{Problem Formulation and Objectives}

\subsection{Problem statement}
Given a stream of physiological samples from a small set of low-cost sensors, produce an honest,
interpretable, real-time estimate of a person's stress that (i) adapts to individual physiology,
(ii) recognizes and abstains under untrustworthy signal conditions, and (iii) never over-claims
beyond what the sensors can support.

\subsection{Objectives}
I defined the project against seven concrete objectives:

\begin{enumerate}[leftmargin=1.4em,itemsep=0.1em]
  \item Build a complete software-first pipeline (data $\rightarrow$ features $\rightarrow$ model
        $\rightarrow$ runtime $\rightarrow$ interface) that runs from replayed or streamed data with
        no hardware dependency.
  \item Detect stress in a personalized manner, calibrated to each user's own baseline.
  \item Match or exceed the $\sim\!0.874$ LOSO benchmark under honest evaluation.
  \item Produce a graded, interpretable output---a $0$--$100$ index with a plain-language
        explanation of which physiological system drives it.
  \item Make abstention a first-class behaviour under motion or poor signal quality.
  \item Enforce honesty structurally, so no clinical or diagnostic claim can be emitted.
  \item Deliver an interface a non-expert can understand, including the system's own limits.
\end{enumerate}

% ===========================================================================
\section{Project Planning and Timeline}

I decomposed the problem into six dependent concerns---data and contracts, feature engineering,
modelling, runtime, serving, and validation/honesty---and adopted a deliberate \emph{software-first}
strategy: prove the entire pipeline on public datasets and a synthetic sensor stream before
attaching hardware, so that model, backend, and interface are all tested against the exact protocol
the device will later use. A fixed raw-event schema defined up front guarantees that a synthetic
generator, a replayed dataset, and future firmware all speak the same protocol.

Table~\ref{tab:gantt} shows the project timeline as a Gantt chart across ten working weeks; shaded
cells denote active weeks.

\begin{table}[t]
\centering
\caption{Project timeline (Gantt chart). Shaded cells indicate active weeks.}
\label{tab:gantt}
\renewcommand{\arraystretch}{1.25}
\setlength{\tabcolsep}{3pt}
\footnotesize
\begin{tabular}{p{5.4cm} *{10}{>{\centering\arraybackslash}p{0.42cm}}}
\toprule
\textbf{Phase} & W1 & W2 & W3 & W4 & W5 & W6 & W7 & W8 & W9 & W10 \\
\midrule
1. Background research \& scoping
 & \cellcolor{barfill} & \cellcolor{barfill} & & & & & & & & \\
2. Data contracts \& dataset loaders
 & \cellcolor{barfill} & \cellcolor{barfill} & \cellcolor{barfill} & & & & & & & \\
3. Feature pipeline (v1)
 & & \cellcolor{barfill} & \cellcolor{barfill} & & & & & & & \\
4. Baseline model (M1) + LOSO
 & & & \cellcolor{barfill} & \cellcolor{barfill} & & & & & & \\
5. Runtime: calibration, gate, states, explanations
 & & & & \cellcolor{barfill} & \cellcolor{barfill} & \cellcolor{barfill} & & & & \\
6. Backend + dashboard (v1)
 & & & & & \cellcolor{barfill} & \cellcolor{barfill} & \cellcolor{barfill} & & & \\
7. Dataset harmonization + held-out validation
 & & & & & & \cellcolor{barfill} & \cellcolor{barfill} & & & \\
8. Accuracy improvement (features-v2, personalization)
 & & & & & & & \cellcolor{barfill} & \cellcolor{barfill} & & \\
9. Real streaming + interface rebuild
 & & & & & & & & \cellcolor{barfill} & \cellcolor{barfill} & \\
10. Documentation, spec, hardening
 & & & & & & & & & \cellcolor{barfill} & \cellcolor{barfill} \\
\bottomrule
\end{tabular}
\end{table}

The two heaviest phases were the runtime (Phase~5, where personalization and abstention live) and
the accuracy work (Phase~8, which raised the model from $0.831$ to $0.919$). Phase~7 produced an
important negative result---that naive dataset pooling harms accuracy---which I acted on by
abandoning pooling.

% ===========================================================================
\section{Methodology}

\subsection{Windowing}
I slide a $60$-second analysis window over the aligned signals, stepping every $30$ seconds. Each
window yields one fixed-schema feature row. A window overlapping a sensor dropout receives a low
coverage score and is routed to the abstention layer (Section~\ref{sec:abstention}) rather than
silently discarded.

\subsection{Feature extraction}
From four channels I compute $19$ absolute features per window using
NeuroKit2~\cite{makowski2021neurokit} for pulse-peak detection and electrodermal decomposition.
These include time-domain heart-rate-variability statistics, frequency-domain HRV, and a
convex-optimization decomposition of skin conductance into tonic and phasic
components~\cite{greco2016cvxeda}. The full feature set is enumerated in Section~\ref{sec:inputs}.
Two features quantify sympathovagal balance and the low/high-frequency ratio of the inter-beat
interval spectrum, defined as
\begin{equation}
\mathrm{LF/HF} = \frac{\displaystyle\int_{0.04}^{0.15} S(f)\,df}{\displaystyle\int_{0.15}^{0.40} S(f)\,df},
\end{equation}
where $S(f)$ is the power spectral density of the inter-beat-interval series; a higher ratio
indicates greater sympathetic dominance.

\subsection{Personalization}
This is the core of the method. Let $x_{i}$ be a physiological feature for window $i$ of a given
user, and let $\mathcal{B}$ be the set of windows in that user's initial quiet calibration period
(the first $\sim\!300$ seconds of accepted data). I compute the personal mean and standard deviation
\begin{equation}
\mu = \frac{1}{|\mathcal{B}|}\sum_{j\in\mathcal{B}} x_j, \qquad
\sigma = \sqrt{\frac{1}{|\mathcal{B}|}\sum_{j\in\mathcal{B}} (x_j-\mu)^2},
\end{equation}
and append the personalized feature
\begin{equation}
x_i^{(p)} = \frac{x_i - \mu}{\max(\sigma,\;\varepsilon)},
\end{equation}
where $\varepsilon$ is a small floor guarding against division by zero. Each physiological feature
is thus represented \emph{twice}---in absolute units and as a personal deviation---doubling the
effective feature count to $36$. Crucially, the reference window $\mathcal{B}$ is defined in
\emph{seconds of signal}, not a fixed count of windows, so that the calibration performed live by the
deployed system exactly matches the reference used during training, regardless of window step.

\subsection{Model architecture}
I use gradient-boosted decision trees (a histogram-based gradient boosting
classifier~\cite{ke2017lightgbm}), a deliberate choice over deep learning. With only $15$ training
subjects and tabular features, tree ensembles are the appropriate and best-performing tool at this
data scale; they handle missing values natively (a dropped sensor simply becomes a NaN, requiring
no imputation); and they remain interpretable. The model has two heads that share the same $36$
inputs:
\begin{itemize}[leftmargin=1.4em,itemsep=0.1em]
  \item \textbf{Head A---graded stress.} A binary classifier (baseline vs.\ stress) wrapped in
        isotonic probability calibration~\cite{zadrozny2002transforming}, so its output is a genuine
        probability $p$. The graded index is $\mathrm{index}=\lfloor 100\,p \rceil$, with ordinal
        bands \emph{calm} ($p<0.45$), \emph{elevated} ($0.45\le p<0.70$), and \emph{high}
        ($p\ge 0.70$).
  \item \textbf{Head B---affect.} A four-class classifier
        (baseline / stress / amusement / meditation) that separates stress arousal from positive
        arousal, since a raised heart rate during an enjoyable moment is not stress and a single
        binary model cannot distinguish the two.
\end{itemize}
Calibration matters: without it a classifier's raw output is a score, not a probability, and the
$0$--$100$ index would be arbitrary. Isotonic calibration fits a monotonic mapping from scores to
empirical probabilities, making the index interpretable.

\subsection{Abstention}
\label{sec:abstention}
Before the model is consulted, a quality gate inspects motion and signal coverage. If dynamic wrist
motion exceeds a threshold, or coverage / pulse-signal quality fall below thresholds, the model is
\emph{not invoked} and the window is reported as \texttt{motion\_paused} or \texttt{poor\_signal}.
On abstention the graded index, level, and affect state are all null. This is a designed feature: a
photoplethysmographic reading during hand movement reflects motion artefact rather than blood
volume, and refusing to score is strictly preferable to emitting a confident but meaningless number.

\subsection{Evaluation protocol}
All reported metrics use grouped leave-one-subject-out cross-validation: for each fold one subject
is entirely withheld, the model is trained on the remainder, and it is scored only on the withheld
subject. I report balanced accuracy---the mean of per-class recalls---which is robust to class
imbalance:
\begin{equation}
\mathrm{BA} = \tfrac{1}{2}\!\left(\frac{TP}{TP+FN} + \frac{TN}{TN+FP}\right),
\end{equation}
compared against a majority-class dummy baseline of $0.5$.

\subsection{Honesty enforcement}
A regression test fails the build if any user-facing string contains clinical terminology (e.g.\
``panic,'' ``anxiety,'' ``burnout,'' ``diagnosis,'' ``clinical,'' ``medical''), making over-claiming
structurally impossible rather than a matter of discipline.

% ===========================================================================
\section{Datasets}

I used three datasets with strictly separated roles, summarized in Table~\ref{tab:data}.

\begin{table}[t]
\centering
\caption{Datasets and their roles. Only WESAD was used for training.}
\label{tab:data}
\footnotesize
\renewcommand{\arraystretch}{1.2}
\begin{tabular}{p{2.5cm} c p{2.0cm} p{6.0cm}}
\toprule
\textbf{Dataset} & \textbf{Subjects} & \textbf{Role} & \textbf{Notes} \\
\midrule
\rowcolor{rowalt}
WESAD~\cite{schmidt2018wesad} & 15 & Training & Validated laboratory stress protocol. I used only the
wrist-worn channels and discarded the cleaner chest sensors deliberately, since the target device
has no chest strap. $\sim\!919$ labelled windows after quality filtering. \\
Stress-Predict~\cite{iqbal2022stresspredict} & 35 & Held out & Different people, a different lab
stress protocol; used to probe generalization, never trained on. \\
\rowcolor{rowalt}
Nurse Stress~\cite{hosseini2022nurse} & 15 & Held out & Real hospital shifts with sparse self-report
labels---the messiest, most naturalistic data; never trained on. \\
\bottomrule
\end{tabular}
\end{table}

I trained \emph{only} on WESAD. My own experiments confirmed the literature: naively pooling WESAD
with Stress-Predict lowered LOSO accuracy from $0.83$ to $0.62$, so I abandoned pooling and retained
the other two datasets strictly as external validation. WESAD's labels derive from a validated
protocol, providing trustworthy ground truth; no equivalent labelled dataset exists for conditions
such as burnout, which is one reason the system does not attempt to predict them.

% ===========================================================================
\section{Inputs}
\label{sec:inputs}

Each $60$-second window yields \kpi{36 model inputs}: $19$ absolute features plus $17$ personalized
counterparts. The two quality/coverage features are not personalized, as they describe the recording
rather than the person. Table~\ref{tab:features} enumerates the absolute features by channel.

\begin{table}[t]
\centering
\caption{The 19 absolute features per window, by sensing channel. Each physiological feature (all
except \texttt{ppg\_quality} and \texttt{valid\_fraction}) additionally contributes a personalized
$(\cdot)^{(p)}$ counterpart, for 36 inputs total.}
\label{tab:features}
\footnotesize
\renewcommand{\arraystretch}{1.2}
\begin{tabular}{p{3.0cm} c p{8.0cm}}
\toprule
\textbf{Channel} & \textbf{\#} & \textbf{Features} \\
\midrule
\rowcolor{rowalt}
Pulse / PPG & 7 & mean heart rate; two HRV measures (SDNN-like, RMSSD); pulse-signal quality; LF
power; HF power; LF/HF ratio \\
Skin conductance / EDA & 7 & tonic level; slope; count and mean amplitude of skin-conductance
responses; tonic mean; tonic slope; phasic mean (cvxEDA) \\
\rowcolor{rowalt}
Skin temperature & 2 & mean; trend (negative sign in the arousal direction---stress \emph{cools} the
skin) \\
Motion / accelerometry & 2 & dynamic motion energy (RMS); 95th percentile (also the abstention
trigger) \\
\rowcolor{rowalt}
Coverage & 1 & fraction of the window with usable samples \\
\bottomrule
\end{tabular}
\end{table}

Two properties of this input design are worth emphasizing. First, EDA contributes $7$ of the $19$
features---tied with pulse for the largest block---and is the purest arousal signal, yet is the
channel most often absent on low-cost hardware. Second, the thermal features carry a \emph{negative}
sign in the arousal direction: because stress cools the skin, a falling temperature \emph{increases}
the modelled arousal, a subtlety that a naive ``higher is worse'' design would get backwards.

% ===========================================================================
\section{Results}

\subsection{Accuracy}
Table~\ref{tab:acc} reports the final model under grouped LOSO on WESAD.

\begin{table}[t]
\centering
\caption{Final model performance under grouped leave-one-subject-out on WESAD.}
\label{tab:acc}
\footnotesize
\renewcommand{\arraystretch}{1.2}
\begin{tabular}{c p{5.2cm} c c}
\toprule
\textbf{Head} & \textbf{Task} & \textbf{Balanced acc.} & \textbf{Chance} \\
\midrule
\rowcolor{rowalt}
A & Graded stress (binary) & \kpi{0.919} & 0.500 \\
B & Affect (4-class) & \kpi{0.616} & 0.250 \\
\bottomrule
\end{tabular}
\end{table}

Head~A additionally attains a macro-F1 of $0.918$. For context, the strongest comparable wrist-only
LOSO result in the literature is $0.874$~\cite{siirtola2019}; the system exceeds it.

\subsection{Ablation: how the accuracy was reached}
Table~\ref{tab:ablation} traces the model's development. Every row is measured identically (WESAD
LOSO), so the steps are directly comparable.

\begin{table}[t]
\centering
\caption{Development of Head~A. Every row is grouped LOSO on WESAD.}
\label{tab:ablation}
\footnotesize
\renewcommand{\arraystretch}{1.2}
\begin{tabular}{p{9.5cm} c}
\toprule
\textbf{Configuration} & \textbf{Balanced acc.} \\
\midrule
\rowcolor{rowalt}
Baseline: 13 features, logistic regression & 0.831 \\
Naive pooling of WESAD + Stress-Predict \emph{(rejected)} & 0.617 \\
\rowcolor{rowalt}
Drop pooling---train on WESAD only & 0.831 \\
$+$ frequency-domain HRV and cvxEDA features & 0.880 \\
\rowcolor{rowalt}
$+$ per-subject baseline personalization & \kpi{0.919} \\
\bottomrule
\end{tabular}
\end{table}

The two gains that survived scrutiny were the richer feature set ($+0.049$) and personalization
($+0.039$). The rejected pooling configuration is retained in the table because it is instructive:
adding data \emph{lowered} accuracy, exactly as the literature predicts for cross-protocol pooling
without domain adaptation.

\subsection{Honest limits}
I report three limitations directly inside the deployed interface, not merely here.

\begin{itemize}[leftmargin=1.4em,itemsep=0.15em]
  \item \textbf{The error bar is non-trivial.} With only $15$ subjects, the standard error on the
        balanced accuracy is $\pm 0.036$.
  \item \textbf{The model is good on average, not uniformly.} Per-subject balanced accuracy ranges
        from $0.50$ (chance) to $1.00$. I therefore present the index as a personal \emph{trend}
        over time, never as a precise instantaneous measurement.
  \item \textbf{$0.919$ is the ceiling for this data.} After reaching it, I ran a systematic tuning
        sweep over window size, prediction smoothing, seed ensembling, calibration length, and
        boosting hyperparameters. Nothing improved on $0.919$ with statistical significance. In the
        process I caught two errors in my own work---a bug in which a calibration parameter silently
        changed with window size, and a case of selecting the best configuration on the same test
        folds used to report it---and I documented both as a formal negative result. The binding
        constraint is the dataset size, not the model.
\end{itemize}

\subsection{System results}
The full automated test suite (over $300$ tests) passes, as does an end-to-end software acceptance
gate. The backend genuinely streams: I verified $18$ windows arriving one at a time over a live
WebSocket connection, correctly tracing the arc from calm through high stress, then
\texttt{motion\_paused}, then \texttt{poor\_signal}. The interface was verified in a real browser: a
first-time user is presented with a single clear action, watches the reading accumulate live, and
sees the model's abstention states and error bar rendered honestly rather than hidden.

% ===========================================================================
\section{Discussion and Limitations}

The results support a clear thesis: at this data scale, \emph{how} the signal is represented matters
more than model sophistication. Personalization and a physiologically motivated feature set together
account for the entire improvement from $0.831$ to $0.919$, while an extensive hyperparameter search
yielded nothing. This is a useful corrective to the reflex of reaching for larger models; with $15$
subjects, the path to higher accuracy runs through more subjects and better labels, not deeper
networks.

The most consequential limitation is dataset size. Fifteen training subjects impose a $\pm 0.036$
standard error and a wide per-subject spread, and no amount of modelling can escape it. A second
limitation is that Head~B's affect classes derive from a single laboratory protocol, so its
distinction between stress and positive arousal should be read as suggestive rather than definitive.
A third is scope: the system measures \emph{physiological arousal relative to baseline}, which is
what the sensors can support, and deliberately makes no claim about specific affective disorders,
because the peripheral signature of anger, fear, and excitement is largely shared and cannot be
disambiguated from these channels alone.

% ===========================================================================
\section{Hardware and Deployment}

The intended device integrates a MAX30102 optical pulse sensor, a BNO inertial measurement unit
(motion), and an MLX90614 infrared skin-temperature sensor, covering three of the four channels. It
currently \emph{lacks a skin-conductance (EDA) sensor}, which accounts for $7$ of the $19$ features
and---as established above---is the cleanest arousal signal. Adding that sensor is the next hardware
step, and I am quantifying its value by retraining the model without EDA on the same data, so the
decision is driven by a measured accuracy cost rather than intuition.

Because the pipeline is built around a fixed raw-event contract, hardware only ever replaces the
\emph{data source}: features, model, explanations, and interface are identical whether data arrive
from a recording, a synthetic generator, or a live board. For demonstration, a held-out recording is
streamed through the exact path the hardware will use, so that the pipeline validated in software is
the pipeline the device runs.

% ===========================================================================
\section{Conclusion and Future Work}

I set out to convert low-cost physiological signals into an honest, personalized, real-time stress
estimate. The system attains a balanced accuracy of $0.919$ under the strictest available
evaluation---above the published wrist-only benchmark---built on a personalization method that adapts
to each individual. Equally important, it is candid about its limits: it abstains when it cannot see,
it reports its own uncertainty, and it is structurally barred from making clinical claims. The
clearest avenues for future work are (i) adding the missing skin-conductance sensor and measuring its
contribution, (ii) re-scoring the two held-out datasets against the current model to quantify
real-world generalization, and (iii) enlarging the training population, which the evidence indicates
is the only route past the current accuracy ceiling.

% ===========================================================================
\begin{thebibliography}{99}
\footnotesize

\bibitem{schmidt2018wesad}
P.~Schmidt, A.~Reiss, R.~Duerichen, C.~Marberger, and K.~Van Laerhoven,
``Introducing WESAD, a multimodal dataset for wearable stress and affect detection,''
in \emph{Proc.\ 20th ACM Int.\ Conf.\ Multimodal Interaction (ICMI)}, 2018, pp.\ 400--408.

\bibitem{siirtola2019}
P.~Siirtola,
``Continuous stress detection using the sensors of commercial smartwatch,''
in \emph{Adjunct Proc.\ ACM Int.\ Joint Conf.\ Pervasive and Ubiquitous Computing (UbiComp/ISWC)},
2019, pp.\ 1198--1201.

\bibitem{greco2016cvxeda}
A.~Greco, G.~Valenza, A.~Lanata, E.~P.~Scilingo, and L.~Citi,
``cvxEDA: A convex optimization approach to electrodermal activity processing,''
\emph{IEEE Trans.\ Biomedical Engineering}, vol.\ 63, no.\ 4, pp.\ 797--804, 2016.

\bibitem{makowski2021neurokit}
D.~Makowski \emph{et al.},
``NeuroKit2: A Python toolbox for neurophysiological signal processing,''
\emph{Behavior Research Methods}, vol.\ 53, no.\ 4, pp.\ 1689--1696, 2021.

\bibitem{ke2017lightgbm}
G.~Ke \emph{et al.},
``LightGBM: A highly efficient gradient boosting decision tree,''
in \emph{Advances in Neural Information Processing Systems (NeurIPS)}, 2017, pp.\ 3146--3154.

\bibitem{zadrozny2002transforming}
B.~Zadrozny and C.~Elkan,
``Transforming classifier scores into accurate multiclass probability estimates,''
in \emph{Proc.\ 8th ACM SIGKDD Int.\ Conf.\ Knowledge Discovery and Data Mining}, 2002,
pp.\ 694--699.

\bibitem{iqbal2022stresspredict}
T.~Iqbal \emph{et al.},
``Stress-Predict dataset: A wearable physiological signals dataset for stress detection,''
\emph{Data in Brief / Sensors}, 2022.

\bibitem{hosseini2022nurse}
S.~Hosseini \emph{et al.},
``A multimodal sensor dataset for continuous stress detection of nurses in a hospital,''
\emph{Scientific Data}, vol.\ 9, no.\ 255, 2022.

\end{thebibliography}

\end{document}
